\documentclass[11pt]{article}

% Minimal ACL review baseline: keep only the official template packages.
\usepackage[review]{acl}
\usepackage{times}
\usepackage{latexsym}
\usepackage[T1]{fontenc}
\usepackage[utf8]{inputenc}
\usepackage{microtype}
\usepackage{inconsolata}
\usepackage{graphicx}

\begin{document}
\section{Introduction}

Language models are increasingly deployed in settings where the value
of a variable changes over time and the model must use the current
value, not an earlier one.
Multi-turn assistants must respect the most recent user correction;
retrieval-augmented systems must use the latest version of a document;
tool-augmented agents must read the current contents of a variable
after a sequence of writes.
Failures of this skill have been documented at the application level,
but the underlying mechanism, and whether the failure is intrinsic to
current models, is not understood.

The task uses $K$ interleaved keys, each receiving $N$ updates, with
the model asked for either the first (FVQ) or current (CVQ) value of
a target key. Both queries are trivially solvable from context, so
failure indicates a retrieval problem, not a length or memory limit.

We evaluate 20 pretrained models spanning 14 open-weight checkpoints
(270M--9B parameters) and 6 proprietary frontier models.
At a representative cell per model block, FVQ$-$CVQ gaps range from
+6 to +53 percentage points, with the majority exceeding +20 pp.
The gap takes two shapes: primacy (FVQ exceeds CVQ; majority of
cells) and reversal (CVQ exceeds FVQ but both below ceiling; smaller
open-weight models at high load).
When we ask for any intermediate position $k$ of a target key
($1 < k < N$), baseline accuracy drops below 5\% within two positions
of the stream start, indicating a general loss of position-indexed
retrieval rather than a first-vs-current confusion.

The failure is not an architectural limit: a small 0.24\%-parameter
LoRA and a training-free format change both restore accuracy to
93--100\% on harder held-out cells.

\end{document}
